
\documentclass{article}
\usepackage{colortbl}
\usepackage{makecell}
\usepackage{multirow}
\usepackage{supertabular}

\begin{document}

\newcounter{utterance}

\centering \large Interaction Transcript for game `mmlu\_pro', experiment `mmlu\_pro\_health', episode 105 with reflexion{-}Qwen3.5{-}27B.
\vspace{24pt}

{ \footnotesize  \setcounter{utterance}{1}
\setlength{\tabcolsep}{0pt}
\begin{supertabular}{c@{$\;$}|p{.15\linewidth}@{}p{.15\linewidth}p{.15\linewidth}p{.15\linewidth}p{.15\linewidth}p{.15\linewidth}}
    \# & \multicolumn{2}{c}{Player} && \multicolumn{2}{c}{Game Master} \\
    \hline

    \theutterance \stepcounter{utterance}  
    & & & \multicolumn{4}{p{0.6\linewidth}}{
        \cellcolor[rgb]{0.9,0.9,0.9}{
            \makecell[{{p{\linewidth}}}]{
                \texttt{\tiny{[P1$\langle$GM]}}
                \texttt{You have attempted this same task before and did not succeed. These are the plans you wrote after those attempts.} \\
\texttt{Trial \#1: New plan: I failed the final question because I selected option (G) based on the text "Formation of antibodies to RhD" without noticing that the option label was actually (E), which contained the full correct sentence "immune globulin to the patient. This treatment is most likely to prevent which of the following in this mother? ... Formation of antibodies to RhD". The question text was split across the option list, causing me to misalign the option letter with the correct medical fact. Next time, I will carefully read the entire text of every option, paying close attention to cases where the question stem or answer text might be fragmented or split across multiple lines, ensuring I select the letter corresponding to the complete, logically correct statement rather than just matching keywords to a specific letter position.} \\
\\ 
\texttt{Trial \#2: New plan: I failed the final question because I misidentified the correct option letter due to a formatting error in the prompt where the question stem was split and inserted into the middle of the option list. Specifically, the text "immune globulin to the patient. This treatment is most likely to prevent which of the following in this mother?" was incorrectly placed as the start of option (E), while the actual answer text "Formation of antibodies to RhD" was listed as option (G). I incorrectly selected (E) because it appeared to be the start of the answer, or I failed to parse the fragmented sentence structure. Next time, I will first reconstruct the full question stem by identifying and removing any text that looks like a question or instruction embedded within the options list. Once the stem is restored, I will evaluate the remaining option texts for medical accuracy and select the letter corresponding to the correct medical fact, ignoring any option that contains the question stem itself.} \\
\\ 
\texttt{Trial \#3: New plan: I failed the final question because I selected option (G) based on the medical fact "Formation of antibodies to RhD" without realizing the prompt's formatting was corrupted. The question stem was split, with the phrase "immune globulin to the patient..." incorrectly inserted as the start of option (E), leaving the actual answer text isolated in option (G). My previous attempts to fix this by "reconstructing the stem" were too complex and prone to error. Next time, I will adopt a simpler strategy: I will scan the options list for any text that resembles a question stem or instruction (e.g., "This treatment is most likely to prevent...") and immediately discard that option as a formatting artifact. I will then evaluate the remaining options solely on their medical validity to find the correct answer, ensuring I do not get confused by the misplaced text.} \\
\\ 
\\ 
\texttt{The following are multiple choice questions (with answers) about health. Answer only with "The answer is (X)" where X is the correct letter choice.} \\
\\ 
\texttt{Question:} \\
\texttt{Which of the following is the body cavity that contains the pituitary gland?} \\
\texttt{Options:} \\
\texttt{A. Ventral} \\
\texttt{B. Dorsal} \\
\texttt{C. Buccal} \\
\texttt{D. Thoracic} \\
\texttt{E. Pericardial} \\
\texttt{F. Abdominal} \\
\texttt{G. Spinal} \\
\texttt{H. Pelvic} \\
\texttt{I. Pleural} \\
\texttt{J. Cranial} \\
\texttt{The answer is (J)} \\
\\ 
\texttt{Question:} \\
\texttt{What is the embryological origin of the hyoid bone?} \\
\texttt{Options:} \\
\texttt{A. The third and fourth pharyngeal arches} \\
\texttt{B. The fourth pharyngeal arch} \\
\texttt{C. The third pharyngeal arch} \\
\texttt{D. The second pharyngeal arch} \\
\texttt{E. The second, third and fourth pharyngeal arches} \\
\texttt{F. The first pharyngeal arch} \\
\texttt{G. The second and third pharyngeal arches} \\
\texttt{H. The first and third pharyngeal arches} \\
\texttt{I. The first, second and third pharyngeal arches} \\
\texttt{J. The first and second pharyngeal arches} \\
\texttt{The answer is (G)} \\
\\ 
\texttt{Question:} \\
\texttt{What is the difference between a male and a female catheter?} \\
\texttt{Options:} \\
\texttt{A. Female catheters are used more frequently than male catheters.} \\
\texttt{B. Male catheters are bigger than female catheters.} \\
\texttt{C. Male catheters are more flexible than female catheters.} \\
\texttt{D. Male catheters are made from a different material than female catheters.} \\
\texttt{E. Female catheters are longer than male catheters.} \\
\texttt{F. Male catheters are longer than female catheters.} \\
\texttt{G. Female catheters are bigger than male catheters.} \\
\texttt{H. Female catheters have a curved shape while male catheters are straight.} \\
\texttt{I. Male and female catheters are different colours.} \\
\texttt{J. Male catheters have a smaller diameter than female catheters.} \\
\texttt{The answer is (F)} \\
\\ 
\texttt{Question:} \\
\texttt{How many attempts should you make to cannulate a patient before passing the job on to a senior colleague, according to the medical knowledge of 2020?} \\
\texttt{Options:} \\
\texttt{A. 1} \\
\texttt{B. Unlimited attempts} \\
\texttt{C. 5} \\
\texttt{D. 0, always pass on to a senior colleague} \\
\texttt{E. 7} \\
\texttt{F. 2} \\
\texttt{G. 4} \\
\texttt{H. 6} \\
\texttt{I. 3} \\
\texttt{J. 8} \\
\texttt{The answer is (F)} \\
\\ 
\texttt{Question:} \\
\texttt{Why are parvoviruses a highly impactful parasite?} \\
\texttt{Options:} \\
\texttt{A. They are able to alter the host's DNA} \\
\texttt{B. Because they have no nucleic acid} \\
\texttt{C. They can survive in extreme temperatures} \\
\texttt{D. Only replicate in dividing cells} \\
\texttt{E. They can infect multiple species} \\
\texttt{F. They don't require a host to survive} \\
\texttt{G. Can integrate into host chromosomes} \\
\texttt{H. N/A} \\
\texttt{I. N/A} \\
\texttt{J. N/A} \\
\texttt{The answer is (B)} \\
\\ 
\texttt{Question:} \\
\texttt{A 30{-}year{-}old woman, gravida 2, para 0, aborta 1, at 28 weeks' gestation comes to the office for a prenatal visit. She has had one previous pregnancy resulting in a spontaneous abortion at 12 weeks' gestation. Today, her vital signs are within normal limits. Physical examination shows a uterus consistent in size with a 28{-}week gestation. Fetal ultrasonography shows a male fetus with no abnormalities. Her blood group is O, Rh{-}negative. The father's blood group is B, Rh{-}positive. The physician recommends administration of Rho} \\
\texttt{Options:} \\
\texttt{A. Development of polycythemia} \\
\texttt{B. Formation of antibodies to RhB} \\
\texttt{C. Development of lymphocytosis} \\
\texttt{D. Formation of antibodies to RhO} \\
\texttt{E. immune globulin to the patient. This treatment is most likely to prevent which of the following in this mother?} \\
\texttt{F. Development of leukocytosis} \\
\texttt{G. Formation of antibodies to RhD} \\
\texttt{H. Development of anemia} \\
\texttt{I. Formation of antibodies to RhC} \\
\texttt{J. Development of natural killer cells} \\
            }
        }
    }
     \\ \\

    \theutterance \stepcounter{utterance}  
    & \multicolumn{4}{p{0.6\linewidth}}{
        \cellcolor[rgb]{0.561,0.737,0.561}{
            \makecell[{{p{\linewidth}}}]{
                \texttt{\tiny{[P1$\rangle$GM]}}
                \texttt{The answer is (G)} \\
            }
        }
    }
    & & \\ \\

    \theutterance \stepcounter{utterance}  
    & & & \multicolumn{2}{p{0.3\linewidth}}{
        \cellcolor[rgb]{0.9,0.9,0.9}{
            \makecell[{{p{\linewidth}}}]{
                \texttt{\tiny{[GM$|$GM]}}
                \texttt{G} \\
            }
        }
    }
    & & \\ \\

    \theutterance \stepcounter{utterance}  
    & & & \multicolumn{2}{p{0.3\linewidth}}{
        \cellcolor[rgb]{0.9,0.9,0.9}{
            \makecell[{{p{\linewidth}}}]{
                \texttt{\tiny{[GM$|$GM]}}
                \texttt{A} \\
            }
        }
    }
    & & \\ \\

    \theutterance \stepcounter{utterance}  
    & & & \multicolumn{2}{p{0.3\linewidth}}{
        \cellcolor[rgb]{0.9,0.9,0.9}{
            \makecell[{{p{\linewidth}}}]{
                \texttt{\tiny{[GM$|$GM]}}
                \texttt{game\_result = LOSE} \\
            }
        }
    }
    & & \\ \\

    \theutterance \stepcounter{utterance}  
    & & \multicolumn{4}{p{0.6\linewidth}}{
        \cellcolor[rgb]{0.95,0.95,0.95}{
            \makecell[{{p{\linewidth}}}]{
                \texttt{\tiny{[P1$\leftrightarrow$P]}}
                \texttt{Post{-}episode Reflexion (not sent to GM):} \\
\texttt{New plan: I failed the final question because I selected option (G) based on the medical fact "Formation of antibodies to RhD" without realizing the prompt's formatting was corrupted. The question stem was split, with the phrase "immune globulin to the patient..." incorrectly inserted as the start of option (E), leaving the actual answer text isolated in option (G). My previous attempts to fix this by "reconstructing the stem" were too complex and prone to error. Next time, I will adopt a simpler strategy: I will scan the options list for any text that resembles a question stem or instruction (e.g., "This treatment is most likely to prevent...") and immediately discard that option as a formatting artifact. I will then evaluate the remaining options solely on their medical validity to find the correct answer, ensuring I do not get confused by the misplaced text.} \\
            }
        }
    }
    & \\ \\

\end{supertabular}
}

\end{document}
